% ================================================================
% Chapter 11: Competing Senders and the Marketplace of Ideas
% ================================================================

\chapter{Competing Senders and the Marketplace of Ideas}
\label{ch:marketplace}

\epigraph{The ideas of the ruling class are in every epoch the ruling ideas,
i.e.\ the class which is the ruling \emph{material} force of society is at the
same time its ruling \emph{intellectual} force.}{Karl Marx, \textit{The German Ideology} (1846)}

% ----------------------------------------------------------------
\section{Introduction}
% ----------------------------------------------------------------

In every human ecology, multiple idea-systems compete for the same audience.
Missionaries contend with shamans for the interpretive loyalty of
newly contacted peoples; Enlightenment rationalism and Romantic nationalism
vie for Europe's educated classes; Instagram influencers and grandparents
offer competing visions of the good life to the same adolescent.
The \emph{marketplace of ideas}---a metaphor usually left at the level of
metaphor---is, in Ideatic Diffusion Theory, a structure amenable to rigorous
analysis.

This chapter formalizes the marketplace.  We introduce the \texttt{IdeaMarket}
structure, which collects a list of competing \emph{senders} (signals), a
receiver's background \emph{repertoire}, and a payoff function defined in
terms of interpretive depth.  A sender $s_1$ \emph{dominates} sender $s_2$
when $s_1$ produces deeper interpretations for every element of the
receiver's repertoire.  We then prove eighteen theorems that expose the
structural constraints on ideological competition: void never dominates,
self-dominance is trivial, depth chains bound competitive advantage, and
coalition of senders is subadditive---connecting to Gramsci's cultural
hegemony, Berger's sacred canopy, and rational-choice theory of religion.

The key insight is that competitive signalling in ideatic space is not a
free-for-all: the subadditivity of depth imposes hard ceilings on how much
any sender---or any coalition of senders---can achieve.  Cultural hegemony,
in the formal sense, is always bounded above by the audience's accumulated
cultural capital.  This structural result places precise limits on the
efficacy of propaganda, evangelism, and ideological persuasion.

Throughout this chapter, we work in an ideatic space $(\IS, \compose, \void, \depth, \res)$
satisfying the IDT axioms.  Signals and repertoire elements are ideas
$s, r \in \IS$; interpretation is composition; payoff is depth.\footnote{The
interpretation model here is identical to that introduced in Chapter~1:
a receiver with background $r$ encountering signal $s$ produces the
interpretation $r \compose s$.}

\subsection{Historical Foundations of the Marketplace Metaphor}

The ``marketplace of ideas'' is among the most influential metaphors in Western
political and intellectual history, yet its genealogy is rarely traced with
care.  The phrase itself is modern, but the concept it encodes---that
truth emerges from the competitive interplay of rival doctrines---has deep
roots in both liberal political theory and the sociology of religion.

John Stuart Mill's \textit{On Liberty} (1859) provides the canonical
philosophical argument.  Mill's defense of free expression rests not on
any natural right but on an \emph{epistemic} claim: suppression of opinion
is harmful because it deprives society of the opportunity to exchange
error for truth, or to sharpen truth by collision with error.  ``The
peculiar evil of silencing the expression of an opinion is, that it is
robbing the human race,'' Mill writes---a statement that, in our
formalism, asserts that removing a sender from $\mathbf{S}$ can only
reduce the maximum achievable payoff across the repertoire $\mathbf{R}$.
Mill's argument implicitly assumes a kind of monotonicity: more senders,
weakly more total payoff.  Our framework makes this assumption precise
and identifies the conditions under which it holds.

The phrase ``marketplace of ideas'' itself is typically attributed to
Justice Oliver Wendell Holmes Jr.'s celebrated dissent in
\textit{Abrams v.\ United States} (1919), in which he wrote: ``the best
test of truth is the power of the thought to get itself accepted in the
competition of the market.''  Holmes's market metaphor imports the
apparatus of classical economics---competition, selection, equilibrium---into
the domain of belief.  In IDT, we take this importation seriously: the
\texttt{IdeaMarket} structure is a formal market, and dominance is a formal
notion of competitive advantage.  But unlike the neoclassical market, the
ideatic market operates under subadditivity rather than constant returns to
scale, and payoff is bounded by receiver capacity rather than by supply.

The sociology of religion provides a parallel tradition.  Rodney Stark and
William Sims Bainbridge, in \textit{A Theory of Religion} (1987), proposed
a rational-choice model of religious competition: churches are ``firms''
producing supernatural compensators, believers are ``consumers'' selecting
among competing suppliers, and the vigor of religious life is a function of
market structure.  Stark's later work, \textit{The Rise of Christianity}
(1996), argued that Christianity's triumph in the late Roman Empire was not
a miracle but a consequence of superior ``market positioning''---a signal
that resonated across a wider range of receiver backgrounds than its
competitors.  Our formal notion of dominance
(Definition~\ref{def:dominance}) captures precisely this competitive
dynamic: a sender dominates when it achieves higher depth on \emph{every}
element of the repertoire.

Richard Dawkins's concept of the \emph{meme}, introduced in \textit{The
Selfish Gene} (1976), offers another precursor.  Dawkins proposed that
cultural units---memes---replicate, mutate, and compete for survival in a
manner analogous to genes.  Daniel Dennett, in \textit{Darwin's Dangerous
Idea} (1995), extended the memetic framework, arguing that cultural
evolution operates through the differential reproduction of ideas in
minds construed as ecological niches.  In our framework, memes are senders,
minds are receivers with repertoires, and ``survival'' corresponds to
achieving non-zero payoff: a meme ``survives'' in a mind when
$\depth(r \compose s) > 0$ for at least one $r \in \mathbf{R}$.

Antonio Gramsci's theory of cultural hegemony, developed in the
\textit{Quaderni del carcere} (1929--1935), provides the critical
counterpoint to liberal marketplace optimism.  For Gramsci, the marketplace
is not a level playing field but a terrain shaped by power.  The ruling
class achieves \emph{egemonia} not merely by producing superior signals but
by structuring the receivers' repertoires---through education, media, and
institutional socialization---so that its signals \emph{necessarily}
dominate.  In IDT terms, hegemony is achieved not by maximizing
$\depth(s)$ but by shaping $\mathbf{R}$ so that
$\depth(r \compose s_{\text{ruling}}) \ge \depth(r \compose s_{\text{rival}})$
for all $r \in \mathbf{R}$.  The theorems of this chapter apply regardless
of how $\mathbf{R}$ was formed; the critical question---\emph{who shapes
the repertoire?}---lies outside the formal system but is illuminated by it.

Edward Herman and Noam Chomsky's ``propaganda model,'' articulated in
\textit{Manufacturing Consent} (1988), extends Gramsci's insight to mass
media.  In their account, the marketplace of ideas is systematically
distorted by ownership concentration, advertising dependency, and
institutional sourcing.  Our payoff bound
(Theorem~\ref{thm:payoff-bound}) provides a structural complement: even
in a perfectly competitive market, payoff is bounded by audience capacity.
Propaganda's efficacy is limited not only by institutional filters but by
the mathematical ceiling that subadditivity imposes.

The modern attention economy, theorized by Herbert Simon (``a wealth of
information creates a poverty of attention''), Tim Wu (\textit{The Attention
Merchants}, 2016), and others, represents the latest iteration of
marketplace competition.  Social media platforms---Facebook's News Feed,
TikTok's For You page, YouTube's recommendation algorithm---are literal
idea markets, with competing senders (creators, advertisers, political
actors) vying for a finite attention budget.  The receiver's repertoire
$\mathbf{R}$ is finite, and each sender's signal competes for interpretive
engagement with each element of $\mathbf{R}$.  The total payoff across all
senders is bounded by the receiver's total interpretive capacity.  The
formal results of this chapter thus apply directly to the digital
marketplace of the twenty-first century: algorithmic curation selects
senders, but the structural constraints---subadditivity, bounded payoff,
resonance preservation---remain invariant.

% ----------------------------------------------------------------
\section{Core Definitions}
\label{sec:ch11-defs}
% ----------------------------------------------------------------

\begin{definition}[Idea Market]
\label{def:idea-market}
An \emph{idea market} is a tuple $\mathcal{M} = (\mathbf{S}, \mathbf{R})$
where $\mathbf{S} = [s_1, \dots, s_m]$ is a list of \emph{competing senders}
(signals) and $\mathbf{R} = [r_1, \dots, r_n]$ is the receiver's
\emph{background repertoire}.  All elements belong to an ideatic space~$\IS$.
\end{definition}

\begin{lstlisting}[caption={IdeaMarket structure in Lean~4}]
structure IdeaMarket (I : Type*) [IdeaticSpace I] where
  senders    : List I
  repertoire : List I
\end{lstlisting}

\begin{definition}[Interpretation]
\label{def:interpret-ch11}
The \emph{interpretation} of a signal $s$ against a repertoire
$\mathbf{R} = [r_1, \dots, r_n]$ is the list
\[
  \operatorname{interpret}(\mathbf{R}, s)
  \;=\; [r_1 \compose s,\; r_2 \compose s,\; \dots,\; r_n \compose s].
\]
Each entry $r_i \compose s$ is the idea produced when background $r_i$
receives signal $s$.
\end{definition}

\begin{definition}[Total Payoff]
\label{def:total-payoff}
The \emph{total interpretive payoff} of signal $s$ on repertoire $\mathbf{R}$
is
\[
  \operatorname{totalPayoff}(\mathbf{R}, s)
  \;=\; \sum_{i=1}^{n} \depth(r_i \compose s).
\]
This aggregates the depth of every interpretation the signal induces.
\end{definition}

\begin{definition}[Dominance]
\label{def:dominance}
Sender $s_1$ \emph{dominates} sender $s_2$ on repertoire $\mathbf{R}$ when
\[
  \forall\, r \in \mathbf{R},\quad
  \depth(r \compose s_1) \;\ge\; \depth(r \compose s_2).
\]
In anthropological terms, $s_1$ is \emph{more culturally productive} than
$s_2$ for this audience.
\end{definition}

\begin{lstlisting}[caption={Dominance in Lean~4}]
def dominates (s_1 s_2 : I) (rep : List I) : Prop :=
  forall r in rep, depth (compose r s_1) >= depth (compose r s_2)
\end{lstlisting}

The notion of dominance is a pointwise ordering on the payoff function.  It
is deliberately strong: $s_1$ must beat $s_2$ on \emph{every} member of the
audience, not merely on average.  This models \emph{Gramscian} hegemony, in
which the ruling ideology must be accepted across all strata of civil
society, not merely by a majority.\footnote{Gramsci's concept of
\emph{egemonia} in the \textit{Quaderni del carcere} requires precisely
this universality: hegemony is not statistical dominance but categorical
dominance across the entire \emph{blocco storico}.}

% ----------------------------------------------------------------
\section{Void and Dominance}
\label{sec:ch11-void}
% ----------------------------------------------------------------

\begin{theorem}[Void Dominates is Identity (11.1)]
\label{thm:void-dom-identity}
The void signal, when received, simply returns the repertoire unchanged:
\[
  \operatorname{interpret}(\mathbf{R}, \void) \;=\; \mathbf{R}.
\]
\end{theorem}

\begin{proof}[Proof sketch]
By induction on the list $\mathbf{R}$.  For each element $r$, the void
axiom gives $r \compose \void = r$, so
$\operatorname{map}(\lambda\, r.\; r \compose \void, \mathbf{R}) = \mathbf{R}$.
\end{proof}

Silence, the null message, is never the \emph{winning} strategy in a
competitive marketplace---it merely preserves the status quo.  Gramsci's
insight is precisely this: hegemony requires \emph{active} cultural
production, not mere silence.  The void signal is the ideational equivalent
of ideological abdication.

\begin{example}[The Roman Marketplace of Cults]
In the late Roman Republic, the traditional Roman pantheon competed with
imported mystery cults---Isis from Egypt, Mithras from Persia, Cybele from
Phrygia---for the devotion of urban populations. Each cult offered a
different ``signal'' to the same audience. The depth of interpretation
varied with the receiver's background: a Stoic philosopher produced a
deeper interpretation of Mithraic ritual than a dock worker, not because
the ritual differed but because $\depth(r_{\text{philosopher}}) >
\depth(r_{\text{worker}})$. The eventual dominance of Christianity can be
modeled as the emergence of a sender whose signal resonated across the
widest range of receiver backgrounds.\footnote{Rodney Stark,
\textit{The Rise of Christianity} (1996), argues that Christianity's
competitive advantage lay precisely in its broad resonance across social
strata.}

The formal apparatus of this chapter illuminates the mechanism.  Each cult
$c_i \in \mathbf{S}$ had its characteristic payoff profile across the
urban repertoire $\mathbf{R}_{\text{Rome}}$.  Mithraism achieved high depth
among soldiers ($\depth(r_{\text{miles}} \compose s_{\text{Mithra}})$ was
large) but near-zero depth among women, whom it excluded.  The cult of Isis
appealed to merchants and sailors but had limited reach among the
senatorial class.  Christianity's competitive strategy---universal
soteriology, egalitarian ecclesiology, portable ritual---produced moderate
but \emph{uniform} depth across the entire repertoire.  It did not dominate
any single cult on its home turf, but it dominated in the aggregate:
$\operatorname{totalPayoff}(\mathbf{R}_{\text{Rome}}, s_{\text{Christ}}) >
\operatorname{totalPayoff}(\mathbf{R}_{\text{Rome}}, s_{c_i})$ for each
rival $c_i$.  The marketplace selected for breadth, not peak depth---a
structural prediction of our payoff framework.
\end{example}

The Roman example illustrates a general principle that recurs throughout
this chapter: in a diverse marketplace, the sender that achieves moderate
depth across \emph{all} receivers often outperforms the sender that achieves
extreme depth in a narrow niche.  This is a direct consequence of the
structure of the total payoff function, which sums depth across the entire
repertoire.  A signal tuned to a single receiver type can achieve high
local depth but low total payoff; a signal tuned to the entire population
achieves lower local depth but higher total payoff.  The tension between
specialization and universality is thus a structural feature of the
marketplace, not merely a strategic choice.

\begin{theorem}[Self-Dominance is Trivial (11.2)]
\label{thm:self-dom}
Every signal dominates itself:
\[
  \forall\, s \in \IS,\quad \operatorname{dominates}(s, s, \mathbf{R}).
\]
\end{theorem}

\begin{proof}[Proof sketch]
For each $r \in \mathbf{R}$, $\depth(r \compose s) \ge \depth(r \compose s)$
holds by reflexivity of~$\le$.
\end{proof}

Self-consistency is the minimal requirement of any ideological system.
That a religion ``dominates itself'' says nothing about competitive
advantage---it is the zero-information baseline.  Philosophically, this is
Hegel's \emph{Sichselbstgleichheit}: self-identity is the formal precondition
of all determination, never the determination itself.

\begin{theorem}[Void Signal Payoff Equals Repertoire Depth (11.3)]
\label{thm:void-payoff}
The total payoff from the void signal equals the total depth of the repertoire:
\[
  \operatorname{totalPayoff}(\mathbf{R}, \void) \;=\; \operatorname{depthSum}(\mathbf{R}).
\]
\end{theorem}

\begin{proof}[Proof sketch]
By Theorem~\ref{thm:void-dom-identity},
$\operatorname{interpret}(\mathbf{R}, \void) = \mathbf{R}$, so
$\operatorname{depthSum}(\operatorname{interpret}(\mathbf{R}, \void))
= \operatorname{depthSum}(\mathbf{R})$.
\end{proof}

A society's interpretive capacity under silence is exactly its accumulated
cultural depth.  This is the \emph{baseline} against which all competing
signals are measured.  In Bourdieu's terms, the void signal reveals nothing
but the audience's pre-existing \emph{capital culturel}.\footnote{Bourdieu,
\textit{La Distinction} (1979), defines cultural capital as the accumulated
stock of cultural knowledge, skills, and dispositions.  Our
$\operatorname{depthSum}(\mathbf{R})$ is its formal analogue.}

\begin{theorem}[Dominance Reflexivity (11.4)]
\label{thm:dom-refl}
Dominance is reflexive: $\operatorname{dominates}(s, s, \mathbf{R})$ for all $s$
and $\mathbf{R}$.
\end{theorem}

\begin{proof}[Proof sketch]
Immediate from Theorem~\ref{thm:self-dom}.
\end{proof}

In the language of partial orders, dominance is a \emph{preorder} (reflexive
and---as we shall see in later chapters---transitive in an appropriate sense).
This is the formal structure underlying marketplace competition: a preorder on
signals, indexed by audiences.

\begin{remark}[Connection to Attention Economics]
The ideatic marketplace anticipates Herbert Simon's observation that
``a wealth of information creates a poverty of attention.'' In our
framework, the receiver's repertoire $\mathbf{R}$ is finite, and each
sender competes for interpretive engagement with each element of
$\mathbf{R}$. The total payoff across all senders is bounded by the
receiver's total interpretive capacity---a formal version of the
``attention budget'' that dominates contemporary media economics.

Simon's insight, formulated in 1971, has acquired new urgency in the
age of algorithmic content curation.  When a platform's recommendation
engine selects which senders to present to a receiver, it is performing
a constrained optimization over the idea market $\mathcal{M}$: choosing
a subset $\mathbf{S}' \subseteq \mathbf{S}$ to maximize (some proxy for)
total payoff subject to the attention constraint $|\mathbf{S}'| \le k$.
The structural results of this chapter---boundedness, subadditivity,
resonance preservation---apply to this optimization problem and constrain
its solutions.  No algorithm, however sophisticated, can circumvent the
subadditivity ceiling.
\end{remark}

\begin{example}[The Protestant Reformation as Marketplace Competition]
The Protestant Reformation, inaugurated by Luther's \textit{Ninety-Five
Theses} in 1517, is perhaps history's most consequential example of a new
sender entering an established ideatic marketplace.  For a millennium,
the Catholic Church had enjoyed something approaching monopoly status in
Western European Christendom: $\mathbf{S} \approx \{s_{\text{Catholic}}\}$.
Luther, Calvin, Zwingli, and the radical reformers introduced competing
signals $s_{\text{Lutheran}}, s_{\text{Calvinist}}, s_{\text{Anabaptist}}$,
each calibrated to different segments of the receiver repertoire.

Luther's signal was optimized for the German-speaking urban middle class:
vernacular scripture, congregational hymns, the doctrine of the
priesthood of all believers.  For a Wittenberg burgher $r_{\text{burgher}}$,
$\depth(r_{\text{burgher}} \compose s_{\text{Luther}}) \gg
\depth(r_{\text{burgher}} \compose s_{\text{Catholic}})$---the Lutheran
signal resonated more deeply with his existing cultural repertoire
(vernacular literacy, anti-clerical sentiment, commercial independence).
But for a Bavarian peasant embedded in a deeply Catholic milieu,
$r_{\text{peasant}}$, the Catholic signal still dominated.  The
Reformation thus illustrates the \emph{audience-dependence} of dominance:
$s_{\text{Luther}}$ dominated $s_{\text{Catholic}}$ on some repertoires
but not others.  The eventual confessional geography of Europe---the
principle of \textit{cuius regio, eius religio}---reflects the fact that
dominance was regional, not universal.

Calvin's Geneva, by contrast, represents a case of \emph{repertoire
reshaping}: the Consistory systematically restructured receivers'
backgrounds (through catechism, discipline, sumptuary laws) so that
$\depth(r_{\text{Genevan}} \compose s_{\text{Calvin}})$ was maximized.
This is Gramscian hegemony avant la lettre---competitive advantage
achieved by shaping the audience, not just the signal.
\end{example}

% ----------------------------------------------------------------
\section{Depth Bounds on Competition}
\label{sec:ch11-bounds}
% ----------------------------------------------------------------

\begin{theorem}[Hermeneutic Payoff Bound (11.5)]
\label{thm:payoff-bound}
The total payoff of any signal $s$ on repertoire $\mathbf{R}$ is bounded:
\[
  \operatorname{totalPayoff}(\mathbf{R}, s)
  \;\le\; \operatorname{depthSum}(\mathbf{R}) + |\mathbf{R}| \cdot \depth(s).
\]
\end{theorem}

\begin{proof}[Proof sketch]
By induction on $\mathbf{R}$.  At each step, the subadditivity axiom gives
$\depth(r \compose s) \le \depth(r) + \depth(s)$, and the bound accumulates
linearly.
\end{proof}

Even the most profound revelation cannot exceed the audience's capacity plus a
linear contribution from the signal itself.  Mass media reaches many, but depth
per receiver is bounded by pre-existing cultural capital.  This theorem is the
formal backbone of what communications scholars call the
\emph{knowledge gap hypothesis}: the information-rich get richer because their
higher $\operatorname{depthSum}(\mathbf{R})$ allows them to extract more from
every signal.\footnote{Tichenor, Donohue, and Olien, ``Mass Media Flow and
Differential Growth in Knowledge,'' \textit{Public Opinion Quarterly}
34 (1970), 159--170.}

\begin{example}[Algorithmic Competition on Social Media Platforms]
Contemporary social media platforms instantiate the idea market structure
with unprecedented literalness.  Consider TikTok's ``For You'' page.  The
platform maintains an enormous sender pool
$\mathbf{S} = \{s_1, s_2, \dots, s_N\}$ (all available videos), and each
user brings a repertoire $\mathbf{R}_u$ shaped by prior viewing history,
demographic background, and expressed preferences.  The recommendation
algorithm selects a subset $\mathbf{S}' \subset \mathbf{S}$ to present to
the user, implicitly attempting to maximize engagement---a proxy for our
total payoff $\operatorname{totalPayoff}(\mathbf{R}_u, s)$.

The subadditivity bound (Theorem~\ref{thm:payoff-bound}) has an immediate
practical consequence: no matter how sophisticated the algorithm, the total
interpretive value extracted is bounded by the user's existing repertoire
depth plus a linear contribution from the content.  This places a hard
ceiling on ``algorithmic radicalization'' theories: the depth of the
interpretation produced by an extremist signal $s_{\text{radical}}$ is
bounded by $\depth(r) + \depth(s_{\text{radical}})$ for each repertoire
element $r$.  A receiver with a shallow repertoire
($\operatorname{depthSum}(\mathbf{R}_u)$ is small) can be moved by the
signal, but only to a bounded degree.  The algorithm cannot manufacture
depth that neither sender nor receiver possesses.

Furthermore, the resonance preservation theorem
(Theorem~\ref{thm:homogeneous-audience}) explains the ``filter bubble''
phenomenon: users with resonant backgrounds
($r_i \res r_j$ for $r_i, r_j \in \mathbf{R}_u$) produce resonant
interpretations of the same content, reinforcing homogeneity.  The platform
does not create the bubble; it amplifies a structural property of
resonance-preserving composition.
\end{example}

The social media example illustrates how the abstract formal results of this
chapter apply to concrete twenty-first-century institutions.  The marketplace
metaphor, born in Mill's liberalism, refined in Holmes's jurisprudence, and
challenged by Gramsci's critical theory, finds its most literal instantiation
in the algorithmic content platforms of the digital age.  Yet the structural
constraints remain invariant: subadditivity, bounded payoff, resonance
preservation.  These are not historical contingencies but mathematical
necessities inherent in the structure of ideatic space.

\begin{theorem}[Single-Element Repertoire Payoff (11.6)]
\label{thm:singleton-payoff}
For a singleton repertoire $[r]$, the payoff of signal $s$ is exactly:
\[
  \operatorname{totalPayoff}([r], s) = \depth(r \compose s).
\]
\end{theorem}

\begin{proof}[Proof sketch]
Direct unfolding of definitions: $\operatorname{interpret}([r], s) =
[r \compose s]$, and $\operatorname{depthSum}([r \compose s]) =
\depth(r \compose s)$.
\end{proof}

The simplest ``audience''---one mind, one background idea---has its payoff
fully determined by the composition of background and signal.  This is the
irreducible unit of cultural reception, the phenomenological atom of
\emph{Verstehen}.

\begin{theorem}[Empty Repertoire Yields Zero (11.7)]
\label{thm:empty-rep-zero}
If no one is listening, every signal has zero payoff:
$\operatorname{totalPayoff}([], s) = 0$.
\end{theorem}

\begin{proof}[Proof sketch]
The empty list maps to the empty list, whose depth sum is zero.
\end{proof}

A prophet without followers, a book without readers---ideational production
without reception is culturally null.  The tree that falls in an empty forest
produces no interpretive depth.  This is not mysticism but arithmetic: zero
receivers yield zero payoff.

% ----------------------------------------------------------------
\section{Competitive Composition}
\label{sec:ch11-comp}
% ----------------------------------------------------------------

\begin{theorem}[Signal Composition Depth Bound (11.8)]
\label{thm:composed-signal-bound}
Composing two signals $s_1, s_2$ yields a compound signal whose depth on
any repertoire element $r$ is triply bounded:
\[
  \depth(r \compose (s_1 \compose s_2))
  \;\le\; \depth(r) + \depth(s_1) + \depth(s_2).
\]
\end{theorem}

\begin{proof}[Proof sketch]
Apply subadditivity twice: first to $s_1 \compose s_2$, then to
$r \compose (s_1 \compose s_2)$, and combine via transitivity.
\end{proof}

Syncretism---composing competing mythologies into a single system---cannot
create more depth than the sum of parts.  The Hellenistic synthesis of Greek
philosophy and Egyptian religion, the syncretic Christianity of colonial Latin
America, the New Age blending of Eastern and Western spiritualities: all are
structurally constrained by this triple bound.  The ``best of both worlds'' is
an arithmetic ceiling, not an aspirational metaphor.

\begin{theorem}[Void Competitor Adds Nothing, Right (11.9)]
\label{thm:void-comp-right}
Adding void as a second sender to compose with $s$ does not change $s$:
$s \compose \void = s$.
\end{theorem}

\begin{proof}[Proof sketch]
This is the right-identity axiom of the monoid.
\end{proof}

\begin{theorem}[Void Competitor Adds Nothing, Left (11.10)]
\label{thm:void-comp-left}
Void composed on the left with any signal yields the signal unchanged:
$\void \compose s = s$.
\end{theorem}

\begin{proof}[Proof sketch]
This is the left-identity axiom of the monoid.
\end{proof}

A ``silent partner'' in cultural production is no partner at all.  Coalitions
with non-contributors are structurally inert.  This is why purely nominal
alliances---the Holy Roman Empire, which Voltaire quipped was ``neither holy,
nor Roman, nor an empire''---add no interpretive depth.  The void contributor
is formally invisible.

The results of this section have deep implications for understanding
cultural syncretism.  The triple bound on composed signals
(Theorem~\ref{thm:composed-signal-bound}) formalizes a constraint that
anthropologists of religion have long observed informally: syncretic
religious systems---Vodou combining West African \emph{orisha} with
Catholic saints, Sikhism blending Hindu devotionalism with Islamic
monotheism, Bahá'í synthesizing Abrahamic and Zoroastrian themes---never
achieve a depth exceeding the sum of their components.  This is not a
failure of synthesis but a structural inevitability.  The ``sum of parts''
ceiling applies universally, regardless of the creativity or profundity
of the synthesizer.

The void-composition results (Theorems~\ref{thm:void-comp-right}
and~\ref{thm:void-comp-left}) complement the syncretism analysis by
establishing the base case: adding nothing to a cultural system changes
nothing.  This seemingly trivial result has a non-trivial implication for
theories of cultural decline.  When a society's cultural production
approaches void---when its intellectuals, artists, and religious
leaders produce signals of negligible depth---the marketplace degenerates.
The total payoff reduces to the existing repertoire's depth, and no
competitive dynamics operate.  This is the formal analogue of Spengler's
``Zivilisation'' phase, in which exhausted cultures merely repeat existing
forms without generating new depth.\footnote{Oswald Spengler,
\textit{Der Untergang des Abendlandes} (1918--1922).  Spengler's
distinction between \emph{Kultur} (creative cultural production) and
\emph{Zivilisation} (exhausted repetition) maps onto the distinction
between non-void and near-void senders in our framework.}

% ----------------------------------------------------------------
\section{Resonance and Competitive Advantage}
\label{sec:ch11-resonance}
% ----------------------------------------------------------------

\begin{theorem}[Resonant Signals Produce Resonant Interpretations (11.11)]
\label{thm:resonant-competitors}
If two competing senders produce resonant signals, $s_1 \res s_2$, then any
receiver with background $r$ produces resonant interpretations:
\[
  s_1 \res s_2
  \quad\Longrightarrow\quad
  (r \compose s_1) \res (r \compose s_2).
\]
\end{theorem}

\begin{proof}[Proof sketch]
By the left-composition preservation of resonance
(\texttt{resonance\_compose\_left}).
\end{proof}

Competing denominations of the same religion---Catholic vs.\ Orthodox, Sunni
vs.\ Shia---produce \emph{resonant} interpretations in the laity.  The
faithful perceive both as ``basically the same religion'' precisely because the
signals resonate.  This theorem is the formal basis of ecumenism: if the
signals resonate, no receiver can distinguish them at the level of
resonance.\footnote{The ecumenical movement, from the 1910 Edinburgh Conference
onward, rests on exactly this structural claim: beneath doctrinal differences,
the core Christian signals resonate.}

\begin{theorem}[Homogeneous Audience (11.12)]
\label{thm:homogeneous-audience}
If two receivers have resonant backgrounds, $r_1 \res r_2$, and receive the
same signal $s$, their interpretations are resonant:
\[
  r_1 \res r_2
  \quad\Longrightarrow\quad
  (r_1 \compose s) \res (r_2 \compose s).
\]
\end{theorem}

\begin{proof}[Proof sketch]
By right-composition preservation of resonance
(\texttt{resonance\_compose\_right}).
\end{proof}

Members of the same cultural milieu---resonant backgrounds---confronted with
the same competing signal will interpret it similarly.  This is why propaganda
works best within culturally homogeneous populations: the more resonant the
audience, the more uniform the reception.  Conversely, propaganda aimed at a
culturally heterogeneous audience (non-resonant backgrounds) produces
divergent interpretations---the structural limit of mass persuasion in diverse
societies.

\begin{theorem}[Full Cultural Resonance (11.13)]
\label{thm:full-resonance}
Resonant backgrounds plus resonant signals imply resonant interpretations:
\[
  r_1 \res r_2 \;\wedge\; s_1 \res s_2
  \quad\Longrightarrow\quad
  (r_1 \compose s_1) \res (r_2 \compose s_2).
\]
\end{theorem}

\begin{proof}[Proof sketch]
By the full composition-preservation of resonance
(\texttt{resonance\_compose}).
\end{proof}

Allied cultures (resonant backgrounds) exposed to similar myths (resonant
signals) will always arrive at compatible worldviews.  This is the formal
mechanism behind what Redfield called the ``Great Tradition / Little
Tradition'' unity: literate urban elites and illiterate peasants, as long as
their backgrounds resonate and the signals they receive resonate, will
inhabit structurally compatible interpretive universes.\footnote{Robert
Redfield, \textit{Peasant Society and Culture} (1956).  Redfield's
bi-partition of tradition into ``great'' and ``little'' variants assumes
precisely this kind of resonance preservation across social strata.}

\begin{remark}[Open Questions in Competitive Signalling]
Several natural questions arise from the framework developed in this
chapter that remain open for future investigation:

\begin{enumerate}
\item \textbf{Dynamic competition.}  Our model is static: senders and
  repertoire are fixed.  In reality, senders \emph{adapt} their signals in
  response to competitive pressure, and receivers' repertoires evolve
  through exposure.  A dynamic extension would model the co-evolution of
  $\mathbf{S}$ and $\mathbf{R}$ over time, potentially connecting to
  evolutionary game theory and replicator dynamics.

\item \textbf{Coalitional signalling.}  When multiple senders form a
  coalition (an alliance of churches, a media conglomerate, a political
  coalition), the coalition's composite signal is the composition
  $s_1 \compose s_2 \compose \cdots \compose s_k$.  The subadditivity
  bound applies, but optimal coalition formation---which subset of senders
  should ally?---is a combinatorial optimization problem whose complexity
  is not yet characterized in our framework.

\item \textbf{Receiver strategic behavior.}  Our model treats receivers as
  passive: they interpret signals with their existing repertoire but do not
  strategically choose which signals to attend to.  Incorporating receiver
  choice would connect the framework to mechanism design and the economics
  of attention.

\item \textbf{Stochastic payoff.}  In practice, depth may be noisy or
  probabilistic---a receiver may or may not ``get'' a given signal on any
  particular exposure.  A stochastic extension of the payoff function would
  connect to information-theoretic models of cultural transmission.
\end{enumerate}

These questions point toward a richer theory of cultural competition that
integrates the static structural results of this chapter with dynamic,
strategic, and stochastic considerations.
\end{remark}

% ----------------------------------------------------------------
\section{Marketplace Structural Theorems}
\label{sec:ch11-structural}
% ----------------------------------------------------------------

\begin{theorem}[Interpretation Preserves Length (11.14)]
\label{thm:interp-length}
The number of interpretations equals the repertoire size:
$|\operatorname{interpret}(\mathbf{R}, s)| = |\mathbf{R}|$.
\end{theorem}

\begin{proof}[Proof sketch]
$\operatorname{interpret}$ is a \texttt{map}, which preserves list length.
\end{proof}

Cultural contact does not change the \emph{number} of categories a society
has---it transforms each one.  The French anthropological term
\emph{bricolage} (Lévi-Strauss) describes precisely this: the bricoleur
rearranges existing schema rather than adding new ones.  Our theorem
formalizes this conservation law of cultural schema.

The length-preservation result has a subtle but important anthropological
implication.  Cultural contact is often described in terms of
\emph{accretion}---new categories being added to a society's conceptual
inventory.  Our theorem shows that, at the formal level of ideatic
composition, interpretation is \emph{transformative} rather than
\emph{additive}: each element of the repertoire is transformed by the
incoming signal, but no new elements are created.  The apparent addition
of new categories in empirical cultural contact (the introduction of
``democracy'' into a society that lacked the concept, for instance)
must be modeled not as a new element appearing in the interpretation list
but as a pre-existing schema being transformed beyond recognition---the
bricoleur's old tool repurposed for a new function.  This is a strong
claim, and it reflects the limitations of our model as much as the
structure of cultural reality.  A richer model might allow the
interpretation function to produce lists of different length, but the
simplicity of the present framework---and its formal tractability---is
purchased at the price of this conservation law.

\begin{theorem}[Additivity of Repertoires (11.15)]
\label{thm:rep-additivity}
Interpreting a signal with a merged repertoire equals concatenating
interpretations of the sub-repertoires:
\[
  \operatorname{interpret}(\mathbf{R}_1 \mathbin{+\!\!\!+} \mathbf{R}_2, s)
  \;=\;
  \operatorname{interpret}(\mathbf{R}_1, s)
  \mathbin{+\!\!\!+}
  \operatorname{interpret}(\mathbf{R}_2, s).
\]
\end{theorem}

\begin{proof}[Proof sketch]
By the distributivity of \texttt{map} over list concatenation.
\end{proof}

When two groups merge---through marriage alliance, conquest, or
migration---their joint interpretation of any signal equals the union of their
separate interpretations.  Cultural integration is structurally additive.
This additivity result is the formal basis of the ``mosaic'' model of
multicultural societies: each group's interpretive contribution is preserved
intact in the combined repertoire.

\begin{theorem}[Composed Signal Associativity (11.16)]
\label{thm:composed-assoc}
Interpreting a composition of signals through background $r$ is the same
as sequential interpretation:
\[
  r \compose (s_1 \compose s_2)
  \;=\;
  (r \compose s_1) \compose s_2.
\]
\end{theorem}

\begin{proof}[Proof sketch]
By associativity of the monoid operation, applied symmetrically.
\end{proof}

Receiving two messages in sequence---first $s_1$, then $s_2$---is
structurally identical to receiving their composition.  The order of cultural
exposure \emph{telescopes} via associativity: a semester's lectures, received
one by one, produce the same interpretation as receiving their composition
all at once.  This is both the strength and the limitation of the model: it
captures the structural content of sequential exposure but not the
psychological dynamics of surprise, anticipation, or fatigue.

\begin{theorem}[Marginal Depth Contribution Bound (11.17)]
\label{thm:marginal-bound}
Adding a new element $r$ to the repertoire increases total payoff by at most
$\depth(r) + \depth(s)$:
\[
  \operatorname{totalPayoff}(r :: \mathbf{R}, s)
  \;\le\;
  \depth(r) + \depth(s) + \operatorname{totalPayoff}(\mathbf{R}, s).
\]
\end{theorem}

\begin{proof}[Proof sketch]
The new payoff contribution is $\depth(r \compose s) \le \depth(r) + \depth(s)$
by subadditivity.
\end{proof}

Educating one more person---adding one more background to the audience---increases
total societal understanding by at most that person's depth plus the signal's.
Diminishing returns in mass education are structurally guaranteed.  This
theorem resonates with Illich's critique of institutional education: past a
certain point, expanding the student body adds progressively less interpretive
value per capita.\footnote{Ivan Illich, \textit{Deschooling Society} (1971).
Illich's argument that schooling yields diminishing returns finds formal
support in our bound.}

\begin{theorem}[Iterated Signal Depth Bound (11.18)]
\label{thm:iterated-signal}
Composing a signal with itself $n$ times yields depth at most $n \cdot \depth(s)$:
\[
  \depth(\compN{s}{n}) \;\le\; n \cdot \depth(s).
\]
\end{theorem}

\begin{proof}[Proof sketch]
By the depth bound on $\operatorname{composeN}$, which follows from induction
and subadditivity.
\end{proof}

Repetition in ritual---chanting a mantra $n$ times, reciting a creed at every
service, replaying a national anthem daily---has bounded interpretive effect.
The subadditivity of depth means repetition is inherently
\emph{anti-inflationary}: you cannot create unbounded meaning through mere
repetition.  This constrains theories of ritual efficacy that posit cumulative
spiritual deepening through sheer repetition.  The effect of the $n$-th
repetition is bounded by $\depth(s)$, not by $n \cdot \depth(s)$.  The mantra
may be powerful, but it is not infinitely so.

\begin{lstlisting}[caption={Iterated signal depth bound in Lean~4}]
theorem iterated_signal_bound (s : I) (n : ℕ) :
    depth (composeN s n) <= n * depth s :=
  depth_composeN s n
\end{lstlisting}

% ----------------------------------------------------------------
\section{Implications for Media Ecology}
\label{sec:ch11-media}
% ----------------------------------------------------------------

The formal results of this chapter, while derived in the abstract setting of
ideatic space, have direct and substantive implications for the interdisciplinary
field of media ecology---the study of how media environments shape human
perception, cognition, and social organization.  We conclude the main body of
the chapter by drawing out four such implications.

\subsection{The Structural Limits of Propaganda}

The payoff bound (Theorem~\ref{thm:payoff-bound}) and the marginal depth
contribution bound (Theorem~\ref{thm:marginal-bound}) together establish a
hard ceiling on the efficacy of propaganda.  Consider a state actor emitting
a propaganda signal $s_{\text{prop}}$ to a population with repertoire
$\mathbf{R}$.  The total interpretive impact is at most
$\operatorname{depthSum}(\mathbf{R}) + |\mathbf{R}| \cdot
\depth(s_{\text{prop}})$.  This bound is independent of the propagandist's
skill, resources, or institutional power.  It depends only on the audience's
pre-existing depth and the signal's intrinsic complexity.

This result is both sobering and reassuring.  It is sobering because it shows
that audiences with deep repertoires---educated, culturally rich
populations---are capable of extracting \emph{more} depth from propaganda, not
less.  The bound increases with $\operatorname{depthSum}(\mathbf{R})$.  It is
reassuring because it shows that propaganda cannot create depth \emph{ex
nihilo}: the ``big lie,'' repeated indefinitely, is bounded by
$n \cdot \depth(s_{\text{lie}})$ (Theorem~\ref{thm:iterated-signal}), not
by $n^2$ or $e^n$.  Repetition is subadditive.

Jacques Ellul's distinction between ``propaganda of agitation'' and
``propaganda of integration'' finds formal expression here.  Agitation
propaganda ($\depth(s)$ is low, targeting a large $|\mathbf{R}|$) achieves
breadth at the expense of depth.  Integration propaganda ($\depth(s)$ is
high, targeting a resonant $\mathbf{R}$) achieves depth at the expense of
breadth.  Neither can escape the payoff bound.\footnote{Jacques Ellul,
\textit{Propagandes} (1962).}

\subsection{Filter Bubbles and Echo Chambers}

The homogeneous audience theorem (Theorem~\ref{thm:homogeneous-audience})
provides a formal account of the ``filter bubble'' phenomenon identified by
Eli Pariser.  When an algorithmic recommender system selects content for a user
whose repertoire elements are mutually resonant ($r_i \res r_j$ for all
$r_i, r_j \in \mathbf{R}_u$), any signal $s$ produces mutually resonant
interpretations: $(r_i \compose s) \res (r_j \compose s)$.  The user's
interpretive world becomes \emph{more} resonant with each exposure, not less.

The formal mechanism is clear: resonance is preserved under composition
(right context), so a resonant repertoire exposed to any signal remains
resonant.  The ``bubble'' is not an artifact of the algorithm but a
structural consequence of resonance preservation.  The algorithm merely
\emph{accelerates} a dynamics that is inherent in the algebra of
ideatic space.  To ``burst'' a filter bubble, one would need to introduce
repertoire elements that are \emph{not} resonant with the existing
repertoire---a structural intervention, not merely an algorithmic one.

\subsection{The Attention Economy as Bounded Competition}

The modern attention economy, in which content creators compete for a finite
resource (human attention), maps directly onto our formal framework.  The
receiver's attention budget corresponds to the finiteness of $\mathbf{R}$:
a receiver can attend to only finitely many signals, and each interpretation
$r \compose s$ consumes interpretive resources.  The total payoff across all
attended signals is bounded by $\operatorname{depthSum}(\mathbf{R}) +
|\mathbf{R}| \cdot \max_s \depth(s)$.

This bound has implications for platform design.  A platform that presents
$k$ signals to a receiver with repertoire of size $n$ achieves total payoff
at most $n \cdot k$ times the sum of average depths.  Increasing $k$ (showing
more content) does not increase payoff per item; it merely spreads the
receiver's interpretive capacity more thinly.  This is the formal basis of
the ``attention scarcity'' that drives the economics of digital platforms:
the scarce resource is not information (signals are abundant) but
interpretive capacity (repertoire is bounded).

\subsection{Cultural Biodiversity and Marketplace Health}

A marketplace with many diverse senders ($|\mathbf{S}|$ is large, senders are
pairwise non-resonant) offers a richer competitive environment than a
marketplace dominated by a few resonant senders.  The full cultural resonance
theorem (Theorem~\ref{thm:full-resonance}) shows that resonant senders produce
resonant interpretations: if $s_1 \res s_2$, then their interpretations in any
receiver are also resonant.  A marketplace dominated by resonant senders thus
produces a \emph{homogeneous} interpretive landscape---many signals, but all
``saying the same thing'' in the resonance sense.

Cultural biodiversity---the analogue of biological diversity in ecological
theory---requires non-resonant senders.  A healthy marketplace, in our
formalism, is one in which the sender pool $\mathbf{S}$ spans a large portion
of ideatic space, ensuring that different receivers with different repertoires
find signals that resonate with their particular backgrounds.  The marketplace
metaphor thus acquires an ecological dimension: just as a healthy ecosystem
requires species diversity, a healthy idea marketplace requires sender
diversity.  Monopolistic control over the sender pool---whether by a state, a
church, or a platform---reduces cultural biodiversity and, with it, the
potential for deep interpretation across the full range of receiver backgrounds.

% ----------------------------------------------------------------
\section{Conclusion}
\label{sec:ch11-conclusion}
% ----------------------------------------------------------------

The marketplace of ideas, when formalized in IDT, reveals a surprising
structural rigidity.  Competition between senders is not anarchic but
constrained: void never wins, dominance is reflexive, payoff is bounded by
audience capacity, composition is subadditive, and resonant signals produce
resonant interpretations regardless of competitive context.

These results formalize intuitions that social theorists have long articulated
in qualitative terms.  Gramsci's hegemony requires active cultural production
(Theorem~\ref{thm:void-dom-identity}); Bourdieu's cultural capital constrains
reception (Theorem~\ref{thm:payoff-bound}); Redfield's Great/Little Tradition
unity rests on resonance preservation (Theorem~\ref{thm:full-resonance});
Illich's diminishing returns in education follow from marginal bounds
(Theorem~\ref{thm:marginal-bound}).

In the next chapter, we turn from competition to \emph{cooperation}: stable
matching of senders and receivers based on cultural compatibility, formalizing
the two-sided matching structures that underlie kinship, marriage exchange,
and interfaith dialogue.
